% ============================================================================
% 第 2 章 · 总体设计
% ============================================================================

\section{系统总体设计}

第一章把多 Agent 协作场景的工程缺口归结为「真实时观测、跨 CLI 异构、单机本地优先、跨节点可扩展、历史可回放」五条诉求。本章把这五条诉求落到伏羲（fuxi）的总体架构上：自顶向下给出设计目标与约束、整体结构、关键数据流、跨节点扩展与性能预算。本章不展开 Rust 实现，只把架构层面的取舍——为何选择事件驱动而非 RPC 主导、为何把 A2A 与事件总线视为正交、为何编排层在玄女与门客之间硬性插一层而非走对等图——逐项说明。

\subsection{设计目标与约束}
\label{subsec:design-goals}

伏羲的总体设计需化解三组矛盾。第一组「真实时」与「不丢消息」——观测端要求事件抵达延迟尽可能低，审计与回放要求事件不能在拥塞时被悄悄丢弃；第二组「多门客并发」与「文件系统隔离」——多个 CLI 门客同时改一个项目目录会撞 git 锁，而完全独立的目录又会让 \texttt{cargo build} 与 \texttt{node\_modules} 这类大型构建产物每次重生；第三组「跨节点扩展」与「单机本地优先」——分布式调度要求节点间共享调度状态，本地优先又要求单节点离线时仍能完整运行。

第一组通过事件总线的双路径架构解决：事件在每次 \texttt{publish()} 同时进入广播路径（Tokio broadcast\parencite{tokio2024broadcast} 零拷贝扇出）与持久化路径（异步 mpsc 队列由后台 writer 串行落库到 SQLite WAL\parencite{sqlite2024wal}）；writer 拥塞时不阻塞调用方，而是 spawn 后台任务等待，并在堆积超阈值时发出哨兵事件让运维侧感知。该设计借鉴 Kafka 的日志驱动事件流经验\parencite{kreps2011kafka}，但不引入 Kafka 这种重型外部组件。第二组通过 git worktree\parencite{git2024worktree} 多工作树机制与 L1/L2/L3 三层沙箱解决：多个门客并发执行时各自在独立 worktree 上写改动，最终合并回主线；只读场景跳过 worktree 创建直接挂载；持久角色级沙箱跨任务复用构建缓存。第三组通过「单机优先 + 钩子化扩展」解决：系统默认在单节点完整运行，所有持久状态都在本地 SQLite；当用户启动多节点部署时，\texttt{fuxi-cli} 注入 \texttt{DistEnqueuer} 与 \texttt{NodeLoadProvider} 两个 trait 对象到编排层，编排层自身不感知分布式细节。

把功能与非功能需求与设计决策做成可校对清单，列于表~\ref{tab:requirement-design-mapping}。这张映射表强制设计决策对需求负责，避免「为做而做」的过度工程。

\begin{table}[H]
\centering
\caption{功能/非功能需求与设计决策映射}
\label{tab:requirement-design-mapping}
\small
\begin{tabularx}{\textwidth}{p{2.6cm} X p{4.0cm}}
\toprule
\textbf{需求类别} & \textbf{具体诉求} & \textbf{对应设计决策} \\
\midrule
多 Agent 调度 & 玄女把任务派给合适门客；同角色多门客并发；异构 CLI 统一调度 & \autoref{subsec:overall-arch} 玄女—门客分层；\autoref{subsec:dataflow} dispatch pump \\
\addlinespace
真实时观测 & 用户在 TUI/IM/PWA 看到任务进展无轮询；事件 P99 $\le$ 50\,ms & EventBus + Firehose；\autoref{subsec:sla} 公式 (2-4) \\
\addlinespace
单一真相源 & 所有事件 append-only 落 SQLite；崩溃后能完整重放 & 双路径架构；\texttt{ReplayCursor} 接口 \\
\addlinespace
工作区隔离 & 多门客并发不撞 git 锁；构建缓存可复用；只读场景轻量化 & L1/L2/L3 三层沙箱 \\
\addlinespace
人工介入 & 门客遇到歧义可悬挂等待；用户回复后能续跑 & A2A \texttt{input-required} 平台映射 \\
\addlinespace
跨节点扩展 & 跨机部署门客；按节点负载自动选择；用户显式 \texttt{@<node>} 钉死 & \autoref{subsec:cross-node} \texttt{host\_nodes} + auto-pin + 钩子注入 \\
\addlinespace
历史回放 & 任意时间点开始重放；回放与 live tail 无缝衔接 & EventStore + 先订阅后查询 \\
\addlinespace
长期记忆 & Agent 间共享事实与角色技能；扩展独立模块不绑主链 & \texttt{fuxi-memory}/\texttt{fuxi-skills} 横向接入 \\
\addlinespace
轻量部署 & 单二进制；无外部数据库；单机离线可用 & SQLite + Tokio + Axum 自包含栈 \\
\addlinespace
可观测背压 & 拥塞不悄悄丢消息；运维侧能看到滞后趋势 & lag 哨兵事件；公式 (2-4) \\
\bottomrule
\end{tabularx}
\end{table}

除上表外还有两条隐性约束：所有 Agent 间协作必须经编排层显式记录，不允许 Agent 直接 RPC 绕过玄女——这是公理「玄女永远有知情权，无否决权」的工程化表达；观测层只读不写，任何业务状态变化不允许从观测端反向修改主链——这条让观测端可被替换或并存（同一时刻浏览器、TUI、PWA 三端共看一份事件流）而不破坏正确性。

\subsection{总体架构}
\label{subsec:overall-arch}

伏羲采用「核心—通讯—编排—执行—观测」五层架构，每层只依赖更下一层暴露的稳定接口，依赖方向严格自下而上，禁止循环引用。整体如图~\ref{fig:overall-architecture} 所示。

\begin{figure}[H]
  \centering
  \includegraphics[width=0.92\textwidth]{drawio/fig-2-1-overall-architecture.pdf}
  \caption{伏羲系统总体架构（五层 + 横向支撑模块）}
  \label{fig:overall-architecture}
\end{figure}

最底层\textbf{核心层}（\texttt{fuxi-core}）集中定义 \texttt{Agent}、\texttt{Workspace}、\texttt{Event}、\texttt{Task} 四个 trait 与对应值对象，不引入任何外部 IO。\textbf{通讯层}由 \texttt{fuxi-events}（事件总线）与 \texttt{fuxi-a2a}（A2A 协议）共同构成，两者职责互补但语义正交——事件总线描述「系统内部已经发生了什么」，是事实流；A2A 协议描述「Agent 之间应当传递什么」\parencite{a2aproject2025a2a}，是契约语义。这一边界划分是伏羲与 AutoGen\parencite{wu2023autogen}、CAMEL\parencite{li2023camel} 这类把协议与状态合并的系统最显著的差异：在伏羲中 A2A 信封只承担最小契约，所有运行时状态都走事件总线。

\textbf{编排层}（\texttt{fuxi-orchestrator}）是中枢，维护门客注册表 \texttt{Shelf}、把任务派给合适的门客、提供事件 republish 通路让门客私有事件流进入公共总线。该层借鉴 MetaGPT\parencite{hong2023metagpt} 的角色分工但更强调「编排不可绕过」——所有门客之间通讯必须经编排层透明记录。\textbf{执行层}由两个适配器（\texttt{fuxi-agent-cc}、\texttt{fuxi-agent-codex}）与隔离层（\texttt{fuxi-workspace}）构成：cc 适配器走 WebSocket 反连\parencite{fette2011websocket} 长连接复用进程，codex 适配器走 spawn-per-dispatch；隔离层基于 git worktree 提供 L1/L2/L3 三层沙箱。\textbf{观测与服务层}由 \texttt{fuxi-firehose}（事件流四端输出）与 \texttt{fuxi-im}（IM 后端）构成：Firehose 把同一份 EventBus 通过 WebSocket、SSE\parencite{whatwg2024sse}、REST 与终端 TUI 四种形态对外暴露；IM 后端面向 PWA 与桌面客户端提供 REST 任务管理与双向消息流。两者均对编排层只读不写。

横向支撑模块（\texttt{fuxi-memory}、\texttt{fuxi-skills}、\texttt{fuxi-scheduler}）通过事件订阅与显式 API 与编排层交互，自身不参与主链调度。\texttt{fuxi-memory} 提供长期记忆，思路源自 Voyager 持久技能库\parencite{wang2023voyager}；\texttt{fuxi-skills} 是角色玉牒加载器；\texttt{fuxi-scheduler} 提供 cron/once/webhook/文件变化四类触发器，统一通过发出 \texttt{TriggerFired} 事件接入主链。最顶层 \texttt{fuxi-cli} 把所有 crate 装配成可执行二进制 \texttt{fuxi}，内部分成 REPL、daemon、dist、im 等十余个子模块。

13 个 crate 总代码规模约 8.25 万行，每个的形式化边界、职责、依赖范围与具体实现细节在第三章 §3.1 模块布局中以表格形式列出，本章只点出依赖关系约束：若令 $C$ 表示所有 crate 集合，关系 $c_i \to c_j$ 表示 $c_i$ 依赖 $c_j$，则系统约束 $\forall c_i, c_j \in C: (c_i \to c_j) \wedge (c_j \to^* c_i) \Rightarrow \bot$（依赖图无环）。这条由 Cargo 在编译期强制检查，比注释规约或代码审查更可靠。\texttt{fuxi-memory} 的「召回结果回传」需要 orchestrator 持有的事件总线，但若让 memory 反向依赖 orchestrator 会形成循环——本设计采用「trait 在低层、impl 在高层」的反向依赖：在 orchestrator 中定义 \texttt{RecallSink} trait，由顶层 \texttt{fuxi-cli} 提供实现并注入。这是规避循环依赖的标准 Rust 工程惯用法，伏羲在四个跨 crate 钩子上一致采用。

\subsection{关键数据流}
\label{subsec:dataflow}

为把架构层面的静态分层映射到运行时动态行为，本节沿一次完整用户请求展示核心数据流。图~\ref{fig:dataflow} 给出从用户输入到事件多端扇出的核心数据流。

\begin{figure}[H]
  \centering
  \includegraphics[width=0.94\textwidth]{drawio/fig-2-2-dataflow.pdf}
  \caption{用户请求的端到端数据流（玄女接收 $\to$ dispatch $\to$ 门客 $\to$ EventBus $\to$ 多端扇出）}
  \label{fig:dataflow}
\end{figure}

数据流分七步。第一步用户在 IM/PWA/TUI 端发出文本指令通过 WebSocket 或 REST 抵达 \texttt{fuxi-im} 的 \texttt{intervene} handler。第二步 handler 调用 \texttt{Fuxi::intervene(xuannv\_id, text)}，编排层先检查玄女当前 \texttt{ShelfStatus}：Idle 则退化为 \texttt{dispatch} 起一轮新 turn（避免 send\_message 在无 active receiver 时被 drop），Busy 则把消息推入 \texttt{PendingOutbox} 等下次 turn 边界 drain。第三步玄女根据 system prompt 与上下文产出工具调用，或选择把子任务派给某个门客。第四步编排层 spawn 后台 pump task，把私有 receiver 的每个事件补上 \texttt{task\_id}/\texttt{agent\_id} 元数据再 republish 到公共总线。第五步 \texttt{EventBus::publish()} 同时投递到广播路径（broadcast 零拷贝扇出）与持久化路径（mpsc + WAL 串行落库）。第六步所有订阅者从广播路径实时收到事件——Firehose 转发四端、IM 路由到对应任务客户端、记忆抽取器订阅 \texttt{DeliverableProduced} 触发候选事实抽取。第七步任务终态抵达后 pump 的 terminal 检测加 500\,ms 缓冲窗口结束 republish，编排层把门客 \texttt{ShelfStatus} 摊回 Idle。

整条数据流有三处关键设计点。其一\textbf{先订阅后查询}的回放协议：当订阅者通过 \texttt{EventBus::replay(cursor, live\_tail=true)} 请求从历史游标接收事件时必须先 \texttt{subscribe} 拿到 broadcast Receiver，再去 SQLite 拉历史，最后无缝拼接。反过来「先查历史再订阅」会留出窗口，期间产生的事件既未入历史快照也未被订阅者收到。其二\textbf{lag 哨兵与背压可观测性}：当 writer 拥塞超阈值（默认 512 条），\texttt{publish()} 通过 broadcast 路径发出 \texttt{Custom \{ "event\_store\_lagged", pending \}} 哨兵事件。哨兵自身有去重逻辑（同一拥塞窗口只发一次），不会反向加剧拥塞。这把背压从「不可见的死锁」变成「可见的告警」\parencite{lamport1978time}。其三 \texttt{input-required} 在伏羲平台的端到端流转：当门客遇到歧义需要人工答复时，伏羲在 wire 上发出 \texttt{state: "input-required"}\parencite{a2aproject2025a2a}，对应当前 A2A 规范 ProtoJSON 枚举 \texttt{TASK\_STATE\_INPUT\_REQUIRED}。编排层把这一状态映射为 \texttt{ShelfStatus::AwaitingInput}，让整体调度器（包括跨节点 \texttt{NodeLoadProvider}）把「在等用户」的门客从「在工作」剥离。这条端到端流转把规范层语义提升为平台级可观测事件流转。

\subsection{跨节点扩展}
\label{subsec:cross-node}

伏羲默认单节点完整运行，但通过显式扩展点支持多节点。当用户开启分布式模式时 \texttt{fuxi-cli} 在 daemon 启动阶段实例化 \texttt{DistController}，注入到编排层 \texttt{DistEnqueuer} 与 \texttt{NodeLoadProvider} 两个 trait 槽位。架构如图~\ref{fig:cross-node}。

\begin{figure}[H]
  \centering
  \includegraphics[width=0.94\textwidth]{drawio/fig-2-3-cross-node.pdf}
  \caption{跨节点部署架构（home 节点 + worker 节点 + 项目 host\_nodes 路由）}
  \label{fig:cross-node}
\end{figure}

跨节点核心数据是 \texttt{Project} 的 \texttt{host\_nodes} 字段——每个项目可声明只能在哪些节点上跑。当任务关联到某项目时 \texttt{Fuxi::auto\_pin\_from\_project} 读出 \texttt{host\_nodes}，结合各节点当前负载选最合适节点。这条 auto-pin 遵循一条规约：当用户已显式 \texttt{@<node>} 钉死任务时，auto-pin 必须立即返回 \texttt{None}，调用方再加 \texttt{task.pinned\_node.is\_none()} 守卫双保险。这是「自动启发式 vs 用户显式意图」边界的明确化。负载选择函数为：
\begin{equation}
\text{node}(p) = \arg\min_{n \in \mathcal{N}_p} \frac{n.\text{inflight}}{\max(n.\text{max\_concurrency},\ 1)}
\end{equation}
$\max(\cdot, 1)$ 处理 \texttt{max\_concurrency = 0} 边界。这与单节点门客选择函数形式相同——伏羲把同一套饱和度比较逻辑既用在节点层也用在门客层。

跨节点事件流通过分布式控制器双向打通：home 通过 reqwest + JSON-RPC 把任务投递到 worker 的 \texttt{DistEnqueuer}；worker 产生的事件通过同一通道反向 republish 回 home 的 EventBus，所有跨节点事件在 \texttt{EventMeta.source\_node\_id} 字段打标。控制器是嵌入式的——它就是 \texttt{fuxi-cli} 中一个子模块，不需要独立运行时。未来若需更强一致性保障可在控制器内部接入 Raft\parencite{ongaro2014raft}，但当前规模（家庭部署不超过 3 节点）下成本不划算。两条易被忽视的兼容性约束：\texttt{Project} 与 \texttt{Task} 所有新增字段必须 \texttt{\#[serde(default)]}，否则旧版 \texttt{meta.json} 与旧 \texttt{events.db} 中的 task JSON 在 \texttt{registry.list()} 时会全部反序列化失败；worker 节点 pre-spawn 的 \texttt{git fetch} 必须 best-effort，失败只 log warn 不挂断 dispatch。

\subsection{SLA 与性能预算}
\label{subsec:sla}

为让「高性能」从主观印象变成可度量约束，本节给出端到端任务派发延迟与 EventBus republish 滞后的形式化预算。记一次任务从用户输入到结果回传的端到端延迟为 $L_{e2e}$，分解为四分量：
\begin{equation}
L_{e2e} = L_{enq} + L_{disp} + L_{exec} + L_{evt}
\end{equation}
其中 $L_{enq}$ 是任务入派发队列开销，$L_{disp}$ 是编排层选门客并发派发指令开销，$L_{exec}$ 是门客实际执行任务耗时（包括模型推理与工具调用），$L_{evt}$ 是事件从产生到抵达观测端的传播开销。在大多数实际工作负载中 $L_{exec}$ 由模型推理时间主导（数秒至数十秒）。本设计为系统层分量设置 P99 预算：$L_{enq}^{P99} \le 5\,\text{ms}$、$L_{disp}^{P99} \le 10\,\text{ms}$、$L_{evt}^{P99} \le 50\,\text{ms}$。$L_{exec}$ 因受模型推理影响不在系统层 SLA 范围。

事件总线 republish 滞后下界由式给出：
\begin{equation}
L_{evt} \geq \max_{i \in \mathcal{S}} \left( T_{\text{serialize}} + T_{\text{broadcast}} + T_{\text{consume},i} \right)
\end{equation}
$\mathcal{S}$ 是订阅者集合，$T_{\text{serialize}}$ 是事件 JSON 序列化耗时，$T_{\text{broadcast}}$ 是 channel 内部分发耗时，$T_{\text{consume},i}$ 是第 $i$ 个订阅者消费耗时\parencite{lamport1978time}。该不等式表明端到端延迟受最慢订阅者制约。伏羲的应对是：广播路径设置 lag 容忍阈值（默认 512 条），超过则丢弃旧事件并发出 \texttt{RecvError::Lagged}——业务侧全量历史走 replay 路径补齐，broadcast 不为慢消费者承担拖累快消费者的代价。

吞吐量与调度损耗形式化为：
\begin{equation}
T = \frac{N}{W_{total}}, \qquad \eta = 1 - \frac{T}{T_{ideal}}, \quad T_{ideal} = \frac{K}{\tau}
\end{equation}
$N$ 是测量窗口内完成的任务数，$W_{total}$ 是该窗口总耗时，$K$ 是门客数，$\tau$ 是单任务理想耗时。$\eta$ 反映通讯链路、调度逻辑与同步开销在整体吞吐中的占比\parencite{kreps2011kafka}。第四章 §4.2 吞吐实验以 $\eta < 0.15$ 作为 1$\sim$8 门客并发场景下的设计目标。公式 (2-2)$\sim$(2-4) 共同给出本文对「高性能」的可度量定义，第三章给每个分量在模块内部的具体优化，第四章给实测数据验证预算是否被满足。

\subsection{本章小结}

本章自顶向下给出了伏羲的总体设计：明确三组核心矛盾的解法并把功能/非功能需求与设计决策做成映射表；给出五层架构与横向支撑模块的整体布局，强调编排层与协议层正交分离让 A2A 保持轻量、观测层只读不写让多端订阅可以并存；沿一次完整用户请求揭示「先订阅后查询」「lag 哨兵」「\texttt{input-required} 端到端流转」三条关键设计点；给出跨节点扩展的钩子化设计与节点选择函数；给出端到端延迟分解、广播滞后下界与吞吐损耗三组形式化公式。下一章基于本章总体设计沿模块依赖自下而上展开核心模块的内部设计与实现。
